\subsection{Defining Rational Exponents: 1/q} 
So far, we have defined all \makeblank exponents.
\begin{definition}
For positive integers $n$ and any real number $x$,
$$x^n=\underbrace{x\cdot x\cdots x}_{n\ \mathrm{copies}}$$
For any real number $x\neq 0$,
$$x^0=1$$
For any negative integer $-n$ and real number $x\neq 0$,
$$x^{-n}=\dfrac{1}{x^n}$$
\end{definition}
Now we'll define $x^{1/q}$ for rational numbers $\frac{1}{q}$. We want this definition to be consistent with the power of a power rule.
\begin{fact}[Power of a Power Rule (integers)]
For any real number $x$ for which all expressions are defined and any integers $n$ and $p$,
$$\left(x^n\right)^p=x^{np}$$
\end{fact}
We define $x^{1/q}$ to be consistent with this definition. In particular:
\begin{lpic}{}{1}
\regtextleft{0.25}{-1}{$\left(x^{1/q}\right)^q={}$};
\end{lpic}
Since $\left(x^{1/q}\right)^q=\makeblanksmall$, it must be that $x^{1/q}=\makeblank[0.75in].$
\begin{definition}
For any positive integer $q$ and any real number $x$ for which $\sqrt[q]{x}$ is defined,
$$x^{1/q}=\sqrt[q]{x}$$
\end{definition}